\section{Discussion}
\label{sec:discussion}

The present study separates two decisions that are usually collapsed in
feature-based algorithm selection: whether to execute an evaluated fixed
landscape query and, conditional on execution, which portfolio action to use.
This formulation is methodological and does not by itself establish that the
query should be skipped or executed in any empirical regime.

\subsection{When is the evaluated query worth executing?}

\TBD{TBD-DISC-01} The statewise utility distribution has not yet been
estimated.  Its eventual interpretation is conditional on the evaluated query,
utility definition, effect direction, interval, and state coverage, and cannot
be extended from the fixed low-cost query to landscape analysis in general.

\subsection{Information in algorithm-agnostic search behavior}

\TBD{TBD-DISC-02} The RQ2 comparison between T0 and the behavior groups remains
undetermined pending nested BBOB-train family-OOF estimates and frozen
evaluation.  Decision utility, AUROC, Average Precision, Spearman correlation,
and Ridge RMSE support different aspects of the comparison; disagreement among
them precludes a single undifferentiated predictive claim.

\subsection{End-to-end resource and performance trade-offs}

\TBD{TBD-DISC-03} The practical resource--performance trade-off is
undetermined.  Its interpretation requires joint evidence from equal-total-FE
final performance, measured query and inference time, query-call rate, and the
cost of calls with non-positive observed utility.  Fewer calls alone would not
establish better resource-aware optimization.

\subsection{Held-out families and cross-benchmark scope}

\TBD{TBD-DISC-04} Effects on held-out BBOB, CEC2017, CEC2022, and engineering
problems remain unestimated and provide no current evidence of generalization.
Any eventual interpretation is suite-specific; query-extraction failures,
divergent effect directions, or incomplete coverage narrow its scope.

\subsection{Search maturity and behavior interpretation}

\TBD{TBD-DISC-05} The RQ5 ablation, coefficient or discriminant-direction
stability, and maturity--utility association remain unestimated.  Coefficients, discriminant
directions, and correlations have a predictive rather than causal
interpretation; a non-monotonic relationship would require support from its
prespecified fit, interval, and stability analysis.

\subsection{Limitations and validity boundaries}

The primary estimand is query-specific: it concerns
\texttt{descriptor\_cheap}, a 16-dimensional custom low-cost descriptor query.
The prespecified standard and broad \texttt{pflacco} queries can assess
configuration robustness but cannot justify claims
about untested landscape representations. The Selection Reference is a fixed
downstream component rather than a claimed algorithm-selection contribution,
and its selector regret must remain visible when interpreting query utility.
The main Decision Model is trained offline on a BBOB function-family split; the
study does not evaluate online controller training as a primary design.

\TBD{TBD-DISC-06} The empirical limitations remain to be quantified through
trajectory and action-loss coverage, query-extraction failures, family-level
dependence, population-transfer frequency, and sensitivity to the three
prespecified query configurations and utility-cost settings.
